\section*{Chapitre \ref{ch:b3:microbiomes} --- Microbiotes et symbioses}

\begin{solution}{exo:b3:microbiomes:1}
\omterm{def:b3:microbiomes:symbiosis}{Mutualisme}~: les deux y gagnent --- légumineuse et rhizobium, corail et
algue. \omterm{def:b3:microbiomes:symbiosis}{Commensalisme}~: l'un y gagne, l'autre est indifférent --- la
plupart des bactéries intestinales qui vivent de ce que l'hôte ne sait
pas digérer. Parasitisme~: l'un y gagne aux dépens de l'autre ---
\emph{Wolbachia} qui tue les embryons mâles. Un même couple peut
basculer~: un commensal intestinal qui franchit la paroi devient un
pathogène, un rhizobium qui cesse de fixer devient un parasite du
nodule, et une algue soumise au stress thermique nuit au corail qui
l'hébergeait.
\end{solution}

\begin{solution}{exo:b3:microbiomes:2}
$H = -(0.5\ln 0.5 + 0.3\ln 0.3 + 0.1\ln 0.1 + 2\times 0.05\ln 0.05) =
0.347 + 0.361 + 0.230 + 0.300 = 1.24$~; $e^{H} = 3.4$ espèces
effectives. $D = 0.25 + 0.09 + 0.01 + 0.0025 + 0.0025 = 0.355$~; $1/D =
2.8$.
\end{solution}

\begin{solution}{exo:b3:microbiomes:3}
L'enquête 16S rapporte qui est là, à peu près au genre, et le nombre
relatif de lectures~; elle ne dit rien des gènes portés, manque les
champignons et les \omterm{def:b3:virology:virus}{virus} (d'autres \omterm{def:b3:genetic-engineering:pcr}{amorces}), et est biaisée par le
nombre de copies et par les \omterm{def:b3:genetic-engineering:pcr}{amorces}. La métagénomique shotgun rapporte
des gènes, des fonctions et, à profondeur suffisante, des génomes et
des souches entières~; elle coûte plus cher par échantillon, est
submergée par l'ADN de l'hôte dans les prélèvements de tissu, et ne
peut toujours pas dire quels gènes sont exprimés ni quelles cellules
sont vivantes.
\end{solution}

\begin{solution}{exo:b3:microbiomes:4}
Fermenter les fibres en \omterm{def:b3:microbiomes:gut}{acides gras à chaîne courte} qui nourrissent le
côlon et l'hôte~; synthétiser la vitamine~K et des vitamines B~;
transformer les acides biliaires et les médicaments~; la \omterm{def:b3:microbiomes:gut}{résistance à
la colonisation} face aux pathogènes~; éduquer le système immunitaire.
Perturbation~: la colite à \emph{Clostridioides difficile} après des
\omterm{def:b3:bacteriology:antibiotics}{antibiotiques} (sont aussi mises en cause~: les maladies inflammatoires
de l'intestin, l'obésité).
\end{solution}

\begin{solution}{exo:b3:microbiomes:5}
$\binom{20}{5} = \num{15504}$. Espèce à $10$~: $1 - \binom{10}{5}/
\num{15504} = 1 - 252/\num{15504} = 0.984$~; à $6$~: $1 - 2002/\num{15504}
= 0.871$~; à $3$~: $1 - 6188/\num{15504} = 0.601$~; à $1$~: $1 -
\num{11628}/\num{15504} = 0.25$. $E[S_{5}] = 2.7$ espèces. Couverture~:
un singleton, $1 - 1/20 = 0.95$.
\end{solution}

\begin{solution}{exo:b3:microbiomes:6}
$400\times 10^{11} = 4\times 10^{13}$ bactéries, soit environ
\qty{40}{g}. Gènes~: $1000$ espèces $\times$ \num{4000}, moins le
recouvrement --- de l'ordre de $3\times 10^{6}$ gènes distincts, plus de
cent fois les \num{20000} de l'homme.
\end{solution}

\begin{solution}{exo:b3:microbiomes:7}
La \omterm{def:b3:microbiomes:nodule}{nitrogénase} est détruite par le dioxygène, mais les bactéroïdes
doivent respirer pour fabriquer seize ATP par N$_{2}$. Le cortex du
nodule forme une barrière de diffusion qui limite l'entrée du
dioxygène, et la \omterm{def:b3:microbiomes:nodule}{leghémoglobine}, à concentration millimolaire, fixe le
dioxygène qui entre de sorte que la concentration libre reste voisine
de \qty{10}{nM} tandis que le flux vers les bactéroïdes reste élevé. Le
rose est celui de l'oxyleghémoglobine~; un nodule vert a dégradé sa
\omterm{def:b3:microbiomes:nodule}{leghémoglobine} --- il est sénescent et ne fixe plus.
\end{solution}

\begin{solution}{exo:b3:microbiomes:8}
Les deux échouent~: pas de nodules et pas de \omterm{def:b3:microbiomes:mycorrhiza}{mycorhize arbusculaire},
parce que l'oscillation calcique est le signal partagé de la \omterm{def:b3:microbiomes:mycorrhiza}{voie
commune de symbiose}. La voie précède donc la nodulation --- elle a été
bâtie pour la \omterm{def:b3:microbiomes:symbiosis}{symbiose} fongique il y a quatre cents millions d'années
et recrutée plus tard par les légumineuses pour les \omterm{def:b3:microbiomes:nodule}{rhizobiums}.
\end{solution}

\begin{solution}{exo:b3:microbiomes:9}
Kiers a exposé certains nodules de soja à une atmosphère d'argon et de
dioxygène, où les \omterm{def:b3:microbiomes:nodule}{rhizobiums} ne pouvaient pas fixer l'azote, tandis que
les autres nodules de la plante fixaient normalement~; la plante a
coupé l'apport de dioxygène aux nodules non fixateurs et leurs
\omterm{def:b3:microbiomes:nodule}{rhizobiums} se sont reproduits environ deux fois moins. Sans une telle
discrimination, des \omterm{def:b3:microbiomes:nodule}{rhizobiums} qui sauteraient la fixation coûteuse se
reproduiraient plus vite que les fixateurs, se répandraient dans la
population, et les plantes finiraient par héberger des bactéries qui ne
donnent rien~: le \omterm{def:b3:microbiomes:symbiosis}{mutualisme} s'effondrerait.
\end{solution}

\begin{solution}{exo:b3:microbiomes:10}
$\binom{N-N_{i}}{n}/\binom{N}{n} = \prod_{j=0}^{n-1}(N - N_{i} - j)/(N -
j)$, chaque facteur étant inférieur à $1$~; passer de $n$ à $n + 1$
multiplie par un facteur de plus inférieur à $1$, de sorte que la
probabilité de manquer l'espèce baisse et que celle de la voir croît
avec $n$~; en $n = N$ le numérateur $\binom{N -
N_{i}}{N}$ est nul et toutes les espèces sont présentes, ce qui donne
$S$. Pour $N_{i},
n \ll N$ chaque facteur vaut environ $1 - N_{i}/N$, le produit environ
$(1 -
N_{i}/N)^{n} \approx e^{-nN_{i}/N} \approx (1 - n/N)^{N_{i}}$ ---
l'approximation du tirage avec remise.
\end{solution}

\begin{solution}{exo:b3:microbiomes:11}
Un mâle est une impasse pour un symbionte qui ne voyage que dans les
œufs, de sorte que tout caractère qui change les mâles en femelles, les
tue au profit de leurs sœurs, ou handicape les femelles non infectées
augmente la transmission du symbionte. L'incompatibilité
cytoplasmique~: le sperme des mâles infectés est modifié de sorte que le
zygote meurt à moins que l'œuf ne porte la même \emph{Wolbachia}, qui
le sauve. Les femelles non infectées perdent donc la descendance
qu'elles conçoivent avec des mâles infectés, tandis que les femelles
infectées n'en perdent aucune~; l'avantage croît avec la fréquence de
l'infection, de sorte qu'au-delà d'un seuil l'infection se fixe même si
les femelles infectées pondent un peu moins --- le coût est plus que
compensé par la fraction détruite des pontes des femelles non
infectées.
\end{solution}

\begin{solution}{exo:b3:microbiomes:12}
L'endosymbionte passe par un goulot de quelques cellules à chaque
génération de l'hôte et ne recombine jamais avec l'extérieur~: la
dérive fixe des mutations légèrement délétères, délétions comprises, et
il n'existe aucune source de gènes de remplacement (le cliquet de
Muller). Les gènes dont l'hôte fournit les produits, ou qui servent à
une vie libre --- enveloppe, motilité, régulation, l'essentiel de la
réparation --- ne manquent à personne et partent les premiers~; la
machinerie de traduction et les gènes qui fabriquent ce dont l'hôte a
besoin (les acides aminés essentiels chez \emph{Buchnera}) sont gardés
le plus longtemps. Les parents libres ont des populations énormes, de
la recombinaison et des apports de gènes, et la sélection entretient
leur génome. Inverser la tendance exigerait des gènes venus du dehors
--- impossible dans l'isolement --- de sorte que les hôtes remplacent
plutôt un symbionte érodé par un neuf, comme plusieurs lignées
d'insectes l'ont fait.
\end{solution}

\begin{solution}{pb:b3:microbiomes:1}
\textbf{1.} $400\times 10^{11} = 4\times 10^{13}$ bactéries, \qty{40}{g}.
\textbf{2.} Un génome de \num{4000} gènes fait environ \qty{4}{Mb}, soit
\qty{4.4}{fg}~:
$4\times 10^{13}\times 4.4\times 10^{-15} = \qty{0.18}{g}$ d'ADN
bactérien, contre
$3\times 10^{13}\times 6.4\times 10^{-12} = \qty{0.19}{g}$ d'ADN
humain --- à peu près autant.
\textbf{3.} $1000\times 0.7\times 4000 + 0.3\times 4000 \approx 2.8\times
10^{6}$ gènes distincts, quelque $140$ fois ceux du génome humain.
\textbf{4.} $30\times 10 = \qty{300}{mmol}$~; $\times 0.9 = \qty{270}{kJ}$~;
\qty{3}{\%} du régime.
\textbf{5.} Non --- \qty{3}{\%} n'explique pas \qty{30}{\%}. Les souris
sans germes absorbent aussi moins efficacement, ont des hormones
intestinales, un stockage des graisses et une thermogenèse altérés, et
un appétit différent~: le \omterm{def:b3:microbiomes:symbiosis}{microbiote} agit sur la physiologie de l'hôte,
et pas seulement comme fournisseur de carburant.
\textbf{6.} Division compensée par l'évacuation~: une population
stationnaire, chaque cellule remplacée chaque jour. Après une
destruction de \qty{99.9}{\%}, il en reste $4\times 10^{10}$ et dix
doublements rétablissent les nombres en une dizaine de jours --- mais
les survivantes et les plus rapides dominent, la composition est
changée (une \omterm{def:b3:microbiomes:gut}{dysbiose}), et dans l'intervalle un envahisseur tel que
\emph{C.~difficile} peut s'installer.
\textbf{7.} Les espèces mineures se partagent $0.4$ à $177$~:
$p = 0.00226$ chacune. $H
= -(0.3\ln 0.3 + 0.2\ln 0.2 + 0.1\ln 0.1 + 0.4\ln 0.00226) = 0.361 +
0.322 + 0.230 + 2.437 = 3.35$~; $e^{H} \approx 28$.
\textbf{8.} $D = 0.09 + 0.04 + 0.01 + 177\times (0.00226)^{2} = 0.141$~;
$1/D = 7.1$. L'\omterm{def:b3:microbiomes:diversity}{indice de Simpson} élève les abondances au carré, de
sorte que les trois espèces dominantes le décident et que les rares n'y
comptent presque pas~; Shannon donne plus de poids aux espèces rares.
\textbf{9.} $1 - 60/2000 = 0.97$~: environ $30$ lectures sur mille
appartiennent à des espèces non encore vues.
\textbf{10.} La richesse croît avec la taille de l'échantillon tant que
des espèces rares continuent d'apparaître~; \num{10000} lectures
atteignent des espèces que \num{2000} manquent. Raréfier le grand
échantillon à \num{2000} lectures par le théorème (ou par
sous-échantillonnage) --- cela devrait donner environ $180$ s'il s'agit
de la même communauté --- et comparer à profondeur égale.
\textbf{11.} $N_{i} = 1$~: $1 - \binom{1999}{1000}/\binom{2000}{1000} =
1 - (2000 - 1000)/2000 = 0.5$. $N_{i} = 5$~: environ
$1 - 0.5^{5} = 0.97$.
\textbf{12.} $e^{2.0} = 7.4$ espèces effectives~; $D \ge 0.5^{2} = 0.25$,
soit environ $0.27$~: une communauté à une espèce dominante et quelques
dizaines de mineures --- l'intestin effondré, de faible diversité, où
\emph{C.~difficile} s'installe.
\textbf{13.} $150\times 8 = \qty{1200}{kg}$ de sucre par hectare, soit
\qty{10}{\%} des \qty{12}{t} photosynthétisées.
\textbf{14.} $150\,000/28 = 5360$ mol de N$_{2}$~; $16\times 5360 =
\num{85700}$ mol d'ATP.
\textbf{15.} $\num{85700}/30 = 2860$ mol de glucose, \qty{514}{kg} ---
environ \qty{43}{\%} du sucre payé. Le reste bâtit et entretient les
nodules et les bactéroïdes, fournit les huit électrons par N$_{2}$ sous
forme de pouvoir réducteur, et assimile l'ammoniac.
\textbf{16.} $1200\times 16 = \qty{19200}{MJ}$ pour \qty{150}{kg}
d'azote~: \qty{128}{MJ/kg}, presque quatre fois les \qty{35}{MJ/kg} de
l'usine --- mais payés en soleil, au champ, sans carburant.
\textbf{17.} Fixateurs $0.9\times 1 = 0.9$~; tricheurs $0.1\times 0.5 =
0.05$~; les tricheurs font $0.05/0.95 = 5.3\,\%$ --- divisés par deux en
une saison.
\textbf{18.} Tricheurs $0.1\times 1.2^{10} = 0.62$ contre fixateurs
$0.9$~: \qty{41}{\%} après dix saisons, en route vers la domination.
Les sanctions retournent la sélection contre le tricheur~; sans elles
le \omterm{def:b3:microbiomes:symbiosis}{mutualisme} s'érode.
\textbf{19.} $10^{6}\times \qty{20}{pg} = \qty{20}{\micro g}$ de carbone
par cm$^{2}$ et par jour~; \qty{18}{\micro g} pour le corail.
\textbf{20.} Respiration \qty{15}{\micro g}~: un excédent de
\qty{3}{\micro g} par cm$^{2}$ et par jour pour la croissance, la
matrice organique de la calcification, le mucus et les gamètes.
\textbf{21.} Avec \qty{5}{\%} des algues, \qty{0.9}{\micro g} par jour
contre \qty{15}{\micro g} respirés~: un déficit de \qty{14}{\micro g}
par jour~; \qty{300}{\micro g} de réserves durent environ trois
semaines.
\textbf{22.} Quatre semaines à $+2$ sont pires~: les dégâts aux
photosystèmes des algues croissent fortement, non linéairement, avec la
température. La mesure fonctionne néanmoins comme alerte parce que le
blanchissement intègre un stress cumulé et que les deux profils
diffèrent de moins que l'incertitude sur le seuil~; c'est un indice
empirique, étalonné sur des blanchissements observés.
\textbf{23.} Un corail stressé peut se tourner vers des espèces d'algues
plus tolérantes à la chaleur, déjà présentes en faible abondance, ou en
acquérir depuis l'eau après un blanchissement~; les algues tolérantes
donnent moins de carbone, de sorte que le corail croît plus lentement,
et la tolérance gagnée est d'un ou deux degrés --- moins que le
réchauffement prévu --- tandis que le rythme du réchauffement dépasse
l'adaptation propre des coraux.
\textbf{24.} $10^{10}$ cm$^{2}$ $\times$ \qty{20}{\micro g} $=
\qty{200}{kg}$ de carbone par jour, soit environ \qty{73}{t} par an.
\textbf{25.} Environ \qty{3}{\%} du régime venus de la fermentation du
\omterm{def:b3:microbiomes:symbiosis}{microbiote}~; \qty{97}{\%} de couverture pour l'échantillon de
séquençage~; \qty{10}{\%} de la photosynthèse payés pour l'azote de la
culture de haricots.
\end{solution}
